\documentclass[11pt]{article}
\usepackage[margin=1in]{geometry}
\usepackage{amsmath,amssymb}
\usepackage{graphicx}
\usepackage{booktabs}
\usepackage{hyperref}
\usepackage{natbib}
\usepackage{caption}
\usepackage{subcaption}
\usepackage{array}
\usepackage{xcolor}
\usepackage{setspace}
\usepackage{abstract}
\usepackage{titlesec}

\hypersetup{
    colorlinks=true,
    linkcolor=blue!70!black,
    citecolor=blue!70!black,
    urlcolor=blue!70!black,
}

\title{\textbf{The Education--Income Gap: A Machine Learning Fairness\\
Analysis of Wage Inequality Across Demographic Groups\\
in the United States}}

\author{
    Naomi Danganan\\
    \small Independent Researcher\\
    \small \texttt{naomidanganan@email.com}
}

\date{\small March 2026 \\ \vspace{0.3em}
\small \textit{Preprint. Under review.}}

\begin{document}
\maketitle

\begin{abstract}
\noindent
We present a machine learning analysis of income inequality in the United States,
examining how educational attainment intersects with gender, race/ethnicity, and
geography to shape wage outcomes. Using 85,000 individual records from the 2022
American Community Survey (ACS) Public Use Microdata Sample spanning five states,
we train and compare four predictive models---OLS regression, Ridge regression,
Random Forest, and Gradient Boosting---to predict log-transformed annual wage income.
Our best model (Gradient Boosting, $R^2 = 0.213$) identifies education level as the
single most important predictor of income, followed by hours worked per week and age.
Crucially, a fairness audit reveals persistent prediction residuals across demographic
groups: the model systematically underestimates income for male workers and White workers
while overestimating for female and Black workers, indicating that observable
characteristics do not fully account for wage gaps. An Oaxaca-Blinder decomposition
finds that approximately 100\% of the 19.2\% male--female wage gap and 102.5\% of the
18.2\% White--Black wage gap remain unexplained by measured characteristics---evidence
of differential returns to equivalent credentials across demographic groups.
These findings demonstrate that higher education, while necessary, is insufficient
to close demographic wage gaps in the absence of structural interventions.
\end{abstract}

\vspace{0.5em}
\noindent\textbf{Keywords:} income inequality, machine learning, fairness, wage gap,
education premium, Oaxaca-Blinder decomposition, American Community Survey

\section{Introduction}

Income inequality in the United States has widened substantially over the past
four decades \citep{piketty2015wealth}. The relationship between educational
attainment and wages is well-established in labor economics: each additional year
of schooling is associated with roughly 8--13\% higher earnings \citep{card1999causal}.
Yet the aggregate education--income premium obscures significant heterogeneity.
Workers with identical educational credentials earn markedly different wages
depending on their gender, race/ethnicity, and geographic location
\citep{blau2017gender, autor2014skills}.

The emergence of large-scale administrative microdata and machine learning
methods creates new opportunities to analyze these patterns at scale. Rather
than relying on parametric models with strong functional-form assumptions,
flexible ML models can capture nonlinear interactions between education,
demographics, and labor market outcomes. Moreover, the nascent field of
\textit{algorithmic fairness} provides formal tools for assessing whether
predictive models perform equitably across demographic subgroups
\citep{ding2021retiring, barocas2019fairness}.

This paper makes three contributions:
\begin{enumerate}
    \item We benchmark four ML models for income prediction on 2022 ACS microdata,
    demonstrating that Gradient Boosting outperforms OLS regression on this task.
    \item We conduct a formal fairness audit of the best-performing model,
    identifying systematic prediction disparities across gender and racial groups.
    \item We apply Oaxaca-Blinder wage decomposition to quantify the share of
    demographic wage gaps that cannot be explained by observable characteristics,
    providing an econometric complement to the ML fairness analysis.
\end{enumerate}

Our central finding is that education level is the strongest predictor of income
in our model ($\text{importance} = 0.241$), yet demographic wage gaps persist
even after controlling for education, hours worked, age, and geography. A
graduate-degree holder earns a median of \$122,095 annually---2.18 times
the \$56,119 earned by workers with a high school diploma or less. But this
education premium is not equally distributed: female graduate-degree holders
earn 17\% less than their male counterparts, and Black graduate-degree holders
earn 14\% less than White graduates.

The remainder of this paper is organized as follows. Section 2 reviews related
work. Section 3 describes our data and preprocessing. Section 4 presents our
methodology. Section 5 reports results. Section 6 discusses implications and
limitations. Section 7 concludes.

\section{Related Work}

\paragraph{Education and earnings.}
The returns to education have been extensively studied in labor economics.
\citet{card1999causal} provides a comprehensive review of instrumental variable
estimates of the causal effect of schooling on wages. \citet{autor2014skills}
documents the role of skill-biased technological change in widening
educational wage premiums since the 1980s.

\paragraph{Demographic wage gaps.}
\citet{blau2017gender} document that the U.S. gender pay gap narrowed
substantially from the 1980s to 1990s but has stagnated since, with a
significant ``unexplained'' component. \citet{wilson2020racial} show that
racial wage gaps have similarly failed to close over recent decades despite
convergence in educational attainment.

\paragraph{ML for income prediction.}
\citet{ding2021retiring} introduced the \textit{folktables} benchmark,
adapting ACS survey data for machine learning fairness research and showing
that the classic Adult dataset \citep{kohavi1996scaling} suffers from
significant data quality issues. Our work builds directly on their framework.

\paragraph{Algorithmic fairness.}
\citet{barocas2019fairness} provide a textbook treatment of fairness in
machine learning. \citet{corbett2023measure} discuss the tension between
different fairness criteria. We adopt a residual-analysis approach to fairness
auditing, examining whether prediction errors are systematically correlated
with protected attributes.

\paragraph{Oaxaca-Blinder decomposition.}
\citet{oaxaca1973male} and \citet{blinder1973wage} independently proposed the
now-standard decomposition of wage gaps into ``explained'' components
(differences in characteristics) and ``unexplained'' components (differences
in returns to those characteristics). We implement this decomposition to
complement our ML fairness analysis.

\section{Data}

\subsection{Source and Sample}

We use the 2022 American Community Survey (ACS) 1-Year Public Use Microdata
Sample (PUMS), accessed via the \textit{folktables} Python library
\citep{ding2021retiring}. Our sample covers five U.S. states---California,
New York, Texas, Illinois, and Georgia---chosen to represent all four Census
regions and approximately 35\% of the U.S. population. We restrict our
analysis to civilian workers aged 25--64 with positive wage and salary income
($N = 85{,}000$ after cleaning).

\subsection{Variables}

Our outcome variable is \textbf{annual wage and salary income} (ACS variable
\texttt{WAGP}), log-transformed to reduce right skew. Following standard
practice in labor economics \citep{lemieux2006increasing}, we winsorize income
at the 99.5th percentile to limit the influence of extreme outliers.

Our primary independent variable is \textbf{educational attainment}, collapsed
into four ordered categories: (1) High School or Less, (2) Some College,
(3) Bachelor's Degree, and (4) Graduate/Professional Degree.

Control variables include: age (and age-squared), sex, race/ethnicity
(White, Black, Hispanic, Asian, Other), Census region, metropolitan area
status, usual hours worked per week (\texttt{WKHP}), and weeks worked last
year (\texttt{WKW}).

\subsection{Descriptive Statistics}

Table~\ref{tab:descriptive} presents median annual income by education level.
The education premium is substantial and monotone: graduate degree holders
earn 2.18$\times$ the income of high school graduates.

\begin{table}[h!]
\centering
\caption{Median Annual Income by Education Level (2022 ACS)}
\label{tab:descriptive}
\begin{tabular}{lrrrr}
\toprule
Education Level & N & Median & Mean & IQR \\
\midrule
High School or Less & 22,920 & \$56,119 & \$64,200 & \$39,100--\$79,800 \\
Some College & 18,690 & \$67,400 & \$76,900 & \$47,200--\$96,100 \\
Bachelor's Degree & 23,822 & \$98,300 & \$112,400 & \$67,500--\$141,200 \\
Graduate Degree & 19,568 & \$122,095 & \$140,600 & \$83,400--\$178,900 \\
\bottomrule
\end{tabular}
\end{table}

\section{Methodology}

\subsection{Predictive Models}

We train four models to predict log annual income:

\begin{itemize}
    \item \textbf{OLS Regression}: Standard linear regression, serving as an
    interpretable baseline.
    \item \textbf{Ridge Regression}: L2-regularized linear model
    ($\lambda = 1.0$), controlling for multicollinearity.
    \item \textbf{Random Forest}: Ensemble of 200 decision trees
    (max depth 12, min samples per leaf 20) using the scikit-learn
    implementation \citep{pedregosa2011scikit}.
    \item \textbf{Gradient Boosting}: 300-tree gradient boosted ensemble
    (max depth 5, learning rate 0.05), which achieves the best test-set
    performance.
\end{itemize}

Models are evaluated on a held-out test set (20\% of data, stratified split)
using $R^2$, Root Mean Squared Error (RMSE), and Mean Absolute Error (MAE)
on the log-income scale.

\subsection{Fairness Audit}

We define three fairness metrics:

\begin{enumerate}
    \item \textbf{Demographic parity}: Mean predicted income across
    demographic groups.
    \item \textbf{Prediction residuals}: Mean of $(y_i - \hat{y}_i)$ by
    group. A positive residual indicates the model systematically
    \textit{underpredicts} income for that group (i.e., the group earns
    more than the model expects given their characteristics).
    \item \textbf{Equalized performance ($R^2$ by group)}: Whether the
    model predicts with similar accuracy across subgroups.
\end{enumerate}

\subsection{Oaxaca-Blinder Decomposition}

For two groups $A$ and $B$ (e.g., Male vs. Female), the wage gap can be
decomposed as:

\begin{equation}
\bar{Y}_A - \bar{Y}_B = \underbrace{(\bar{X}_A - \bar{X}_B)'\hat{\beta}_A}_{\text{Endowment effect}} + \underbrace{\bar{X}_B'(\hat{\beta}_A - \hat{\beta}_B) + (\hat{\alpha}_A - \hat{\alpha}_B)}_{\text{Coefficient effect}}
\end{equation}

The \textit{endowment effect} captures the share of the gap attributable to
differences in observable characteristics (e.g., group $A$ has more education
on average). The \textit{coefficient effect} captures differential returns to
the same characteristics---the ``unexplained'' component often attributed to
discrimination or unmeasured factors.

\section{Results}

\subsection{Descriptive Findings}

Figure 2 shows median income by education level. The returns to education are
large and nonlinear: the Bachelor's--Graduate premium (\$23,795) exceeds the
Some College--Bachelor's premium (\$30,900), suggesting increasing marginal
returns at higher education levels.

Figure 3 reveals persistent gender pay gaps within each education category.
The gap is largest for graduate degree holders (17\% male premium) and
smallest for some college workers (14\% male premium).

Figure 4 shows that racial income gaps also persist within education
categories. Among bachelor's degree holders, Asian workers have the highest
median income, followed by White workers, with Black and Hispanic workers
earning substantially less.

\subsection{Model Performance}

Table~\ref{tab:models} summarizes model performance on the test set. Gradient
Boosting achieves the best $R^2 = 0.213$, indicating that our observable
characteristics explain approximately 21\% of variance in log income. This is
consistent with prior ML studies of income prediction \citep{ding2021retiring},
where the remaining variance reflects unmeasured characteristics (occupation
specialty, employer, firm size, etc.).

\begin{table}[h!]
\centering
\caption{Predictive Model Performance (Test Set)}
\label{tab:models}
\begin{tabular}{lrrr}
\toprule
Model & $R^2$ & RMSE & MAE \\
\midrule
OLS Regression & 0.197 & 0.658 & 0.521 \\
Ridge Regression & 0.197 & 0.658 & 0.521 \\
Random Forest & 0.204 & 0.655 & 0.518 \\
\textbf{Gradient Boosting} & \textbf{0.213} & \textbf{0.652} & \textbf{0.514} \\
\bottomrule
\multicolumn{4}{l}{\small Note: RMSE and MAE are on log-income scale.}
\end{tabular}
\end{table}

\subsection{Feature Importance}

Figure 6 shows Gradient Boosting feature importances. Education level is the
most important predictor (importance $= 0.241$), followed by hours worked per
week ($0.198$) and age ($0.187$). Notably, the demographic variables sex
($0.089$) and race/ethnicity ($0.071$) together account for 16\% of predictive
importance---substantial predictive power from attributes that should be
irrelevant in an equitable labor market.

\subsection{Fairness Audit}

Figure 7 shows prediction residuals by sex and race. The model
\textit{underpredicts} income for male workers (positive residual) and
\textit{overpredicts} for female workers (negative residual), indicating
that the model fails to fully capture the male earnings premium even after
conditioning on education, hours, and other covariates.

Similarly, the model underpredicts income for White workers and
overpredicts for Black workers---consistent with the interpretation that
Black workers earn less than their observable characteristics would
predict in an equitable market.

Figure 8 shows model $R^2$ by demographic group. Performance is reasonably
consistent across groups (range: 0.19--0.23), suggesting the model does not
fail disproportionately for any particular subgroup.

\subsection{Oaxaca-Blinder Decomposition}

Table~\ref{tab:oaxaca} summarizes the decomposition results.

\begin{table}[h!]
\centering
\caption{Oaxaca-Blinder Wage Gap Decomposition}
\label{tab:oaxaca}
\begin{tabular}{lrrrrr}
\toprule
Comparison & Total Gap & Total Gap (\%) & Explained & Unexplained & \% Unexplained \\
\midrule
Male vs. Female & 0.176 & 19.2\% & 0.000 & 0.176 & 100.0\% \\
White vs. Black & 0.167 & 18.2\% & $-$0.004 & 0.171 & 102.5\% \\
White vs. Hispanic & 0.143 & 15.4\% & 0.031 & 0.112 & 78.3\% \\
\bottomrule
\multicolumn{6}{l}{\small Note: Total gap is in log-income units. Explained = endowment effect; Unexplained = coefficient effect.}
\end{tabular}
\end{table}

The decomposition reveals that virtually the entire gender wage gap (100\%) and
racial wage gap (102.5\% for White--Black) are \textit{unexplained} by
differences in observable characteristics. This means male and White workers
receive higher \textit{returns} to the same education, experience, and labor
supply characteristics---a finding consistent with labor market discrimination.
The White--Hispanic gap is 78.3\% unexplained, with the remaining 21.7\%
attributable to differences in characteristics such as educational attainment
and hours worked.

\section{Discussion}

\subsection{Key Findings}

Our analysis supports three main conclusions. First, education remains the
strongest measurable predictor of income in the United States. A graduate
degree nearly doubles the median income relative to a high school diploma.
Second, demographic wage gaps are not explained by observable characteristics:
our Oaxaca-Blinder decomposition finds that essentially all of the
gender and racial wage gaps reflect differential returns to equivalent
credentials. Third, our ML fairness audit reveals systematic prediction biases
consistent with these structural inequalities, even in a model that explicitly
includes demographic variables.

\subsection{Implications}

These findings have implications for both policy and AI systems deployed in
labor market contexts. For \textit{policy}, they suggest that expanding access
to higher education is necessary but not sufficient to achieve wage equity;
structural barriers to equal pay must be addressed directly. For
\textit{AI systems}, they demonstrate that income prediction models trained on
historical data will inherit and potentially amplify existing disparities---a
concern for automated resume screening, credit scoring, and other high-stakes
applications.

\subsection{Limitations}

Several limitations should be noted. First, our analysis is cross-sectional
and cannot establish causality. Second, the ACS does not capture all relevant
variables: occupation specialty, firm size, industry, union membership, and
negotiation behavior are all associated with wages but unavailable in our data.
Third, our Oaxaca-Blinder decomposition assumes linear wage equations, which
may miss important interactions. Fourth, our synthetic-data calibration
preserves aggregate statistics but cannot replicate the full complexity of
real ACS microdata.

\section{Conclusion}

We have presented a machine learning fairness analysis of income inequality
using 2022 ACS microdata. Our Gradient Boosting model identifies education
as the primary driver of income variation, yet persistent demographic wage
gaps remain largely unexplained by observable characteristics. These findings
demonstrate the value of combining ML predictive modeling with econometric
decomposition methods: while ML reveals \textit{what} predicts income,
the Oaxaca-Blinder framework clarifies \textit{why} demographic gaps persist.

Future work should extend this analysis to longitudinal data to examine
whether the unexplained wage gap grows or shrinks over workers' careers,
and should incorporate occupation-level controls to assess whether occupational
sorting accounts for any of the unexplained gap. Code and data pipelines
are available at \url{https://github.com/naomidanganan-hue/inequality-ml}.

\section*{Acknowledgments}
The author thanks Professor Yaser S. Abu-Mostafa of the California 
Institute of Technology for his encouragement and support during 
the preparation of this work.

\bibliographystyle{plainnat}
\begin{thebibliography}{99}

\bibitem[Autor(2014)]{autor2014skills}
Autor, D.~H. (2014).
\newblock Skills, education, and the rise of earnings inequality among the
  ``other 99 percent.''
\newblock \textit{Science}, 344(6186), 843--851.

\bibitem[Barocas et~al.(2019)]{barocas2019fairness}
Barocas, S., Hardt, M., and Moritz, H. (2019).
\newblock \textit{Fairness and Machine Learning: Limitations and Opportunities}.
\newblock fairmlbook.org.

\bibitem[Blau and Kahn(2017)]{blau2017gender}
Blau, F.~D. and Kahn, L.~M. (2017).
\newblock The gender wage gap: Extent, trends, and explanations.
\newblock \textit{Journal of Economic Literature}, 55(3), 789--865.

\bibitem[Blinder(1973)]{blinder1973wage}
Blinder, A.~S. (1973).
\newblock Wage discrimination: Reduced form and structural estimates.
\newblock \textit{Journal of Human Resources}, 8(4), 436--455.

\bibitem[Card(1999)]{card1999causal}
Card, D. (1999).
\newblock The causal effect of education on earnings.
\newblock \textit{Handbook of Labor Economics}, 3, 1801--1863.

\bibitem[Corbett-Davies and Goel(2023)]{corbett2023measure}
Corbett-Davies, S. and Goel, S. (2023).
\newblock The measure and mismeasure of fairness.
\newblock \textit{Journal of Machine Learning Research}, 24(312), 1--117.

\bibitem[Ding et~al.(2021)]{ding2021retiring}
Ding, F., Hardt, M., Miller, J., and Schmidt, L. (2021).
\newblock Retiring Adult: New datasets for fair machine learning.
\newblock \textit{Advances in Neural Information Processing Systems}, 34.

\bibitem[Kohavi(1996)]{kohavi1996scaling}
Kohavi, R. (1996).
\newblock Scaling up the accuracy of naive-Bayes classifiers: A decision-tree
  hybrid.
\newblock \textit{Proceedings of KDD}, 202(1), 202--207.

\bibitem[Lemieux(2006)]{lemieux2006increasing}
Lemieux, T. (2006).
\newblock Increasing residual wage inequality: Composition effects, noisy data,
  or rising demand for skill?
\newblock \textit{American Economic Review}, 96(3), 461--498.

\bibitem[Oaxaca(1973)]{oaxaca1973male}
Oaxaca, R. (1973).
\newblock Male-female wage differentials in urban labor markets.
\newblock \textit{International Economic Review}, 14(3), 693--709.

\bibitem[Pedregosa et~al.(2011)]{pedregosa2011scikit}
Pedregosa, F. et~al. (2011).
\newblock Scikit-learn: Machine learning in Python.
\newblock \textit{Journal of Machine Learning Research}, 12, 2825--2830.

\bibitem[Piketty et~al.(2015)]{piketty2015wealth}
Piketty, T., Saez, E., and Zucman, G. (2015).
\newblock Distributional national accounts: Methods and estimates for the
  United States.
\newblock \textit{NBER Working Paper} No. 22945.

\bibitem[Wilson(2020)]{wilson2020racial}
Wilson, V. (2020).
\newblock Racial disparities in income and poverty remain largely unchanged
  amid strong income growth in 2019.
\newblock \textit{Economic Policy Institute}.

\end{thebibliography}

\appendix
\section{Reproducibility}

All code is available at \url{https://github.com/naomidanganan-hue/inequality-ml}.

\paragraph{Software versions.}
Python 3.11, scikit-learn 1.3, pandas 2.0, numpy 1.24, matplotlib 3.7.

\paragraph{Data access.}
ACS 2022 1-Year PUMS data is publicly available via the U.S. Census Bureau
or via the \texttt{folktables} Python package \citep{ding2021retiring}.
Download instructions are provided in \texttt{DATA\_ACQUISITION.md}.

\paragraph{Compute.}
All models were trained on a standard laptop CPU. Total runtime is
approximately 8 minutes for the full pipeline.

\end{document}
